\section{Introduction}
\label{sec:introduction}

Deploying large language models (LLMs) for domain-specific quantitative analysis raises a first-order validation problem: distinguishing outputs produced by genuine structural reasoning about the domain from outputs produced by surface-level pattern recall of training-corpus text \citep{mccoy2023embers}. The problem is especially acute in finance, where news and analyst coverage of well-known events (the 2020 COVID crash, the 2021 meme-stock episode, the 2023 banking stress) are heavily represented in training data, and where most structurally interesting patterns occurred in years the model has seen.

This paper proposes \textbf{temporal obfuscation testing} as a validation methodology for this problem and applies it to options dealer gamma-exposure (GEX) regimes. The method strips calendar dates, ticker symbols, and any contextual identifiers from input sequences, and re-evaluates an LLM on the stripped input. When obfuscated detection rates remain high and discriminate persistent regimes from synthetic controls, the model's performance cannot be attributed to memorization of the underlying dates; when they collapse, memorization is the parsimonious explanation. We use dealer GEX as the demonstration domain because it combines mechanically grounded constraints (dealers \emph{must} hedge delta under regulatory and risk-management rules) with a sharp natural contrast---2020 (pre-0DTE, episodic dealer positioning) versus 2024 (post-0DTE, sustained structural dealer positioning)---that a genuinely reasoning system should separate on structural grounds alone \citep{ni2005stock,garleanu2009demand,dim2023odtes}.

\paragraph{Research gap.}
Prior work has, independently, (i)~characterised dealer-gamma hedging and its microstructure effects (\citealp{ni2005stock,garleanu2009demand,anderegg2022impact,dim2023odtes,adams2025zerodte}; among others); (ii)~documented the rapid growth of zero-days-to-expiration (0DTE) options between 2022 and 2025 and its implications for intraday volatility \citep{cboe2024zero,fishman2023gamma,cboe2025spx0dte}; and (iii)~evaluated LLM reasoning with probing and chain-of-thought techniques in non-financial domains \citep{wei2022chain,kojima2022large,mccoy2023embers}. What is missing is a method for validating LLM structural reasoning \emph{within} financial microstructure that (a)~controls for training-data memorisation of specific events and dates, (b)~is empirically tested at a scale comparable to the domain it targets, and (c)~distinguishes genuine structural detection from reproduction of a volatility-regime classifier. Temporal obfuscation testing, in combination with the multi-scale validation design in this paper and the Markov-switching benchmark comparison in \S\ref{sec:regime:benchmark}, is designed precisely for this gap.

\paragraph{Why 0DTE matters here.}
The 2022--2025 growth of 0DTE options in SPY and SPX is a natural setting for an obfuscation study because it created an observable structural shift in dealer positioning within the training horizon of modern LLMs. By 2024, 0DTE volume had risen to approximately 46\% of total SPY options volume and 59\% of SPX volume by 2025 \citep{cboe2024zero,cboe2025spx0dte,dim2023odtes,adams2025zerodte}, concentrating gamma exposure at ultra-short maturities. If an LLM reports 2024 as a persistent-regime year and 2020 as a fragmented-regime year \emph{after} all dates and tickers are stripped, it is detecting a structural property of the numerical sequence, not recalling that 2024 contained the word ``0DTE'' in the training corpus. That is the epistemic move temporal obfuscation is designed to enable.

We validate the framework at two temporal scales, moving from single-day pattern detection to persistent 30-day regime identification---a qualitative strengthening of the validation problem, not merely a change in window length.

\textbf{Single-day validation.} Applied to 242 trading days (SPY, 2024), obfuscation testing achieves 71.5\% detection of dealer hedging patterns using unbiased prompts, with 90.9\% of detections materializing in forward returns. A \textbf{raw chain validation} removing all pre-calculated metrics achieves 92.3\% detection---outperforming the GEX-assisted baseline by 30.8 percentage points---demonstrating that LLMs reconstruct dealer positioning from first principles rather than matching parametric summaries \citep{regan2025obfuscation}.

\textbf{Multi-day regime detection.} Extending to 30-day windows across six years (2020--2025), the framework achieves 81.2\% detection of persistent regimes in 2024 (95\% CI [75.8, 86.1]\%) versus 12.1\% in 2020 (95\% CI [8.1, 16.6]\%) --- a 69.1 percentage point separation, $\varphi = 0.69$, Fisher's exact $p = 1.8 \times 10^{-52}$ --- with 0\% false positives on synthetic controls. Multi-year analysis reveals gradual regime evolution tracking 0DTE adoption: detection rates rise from 3.7\% (2021) to 100\% (2024), with average GEX magnitude growing from \$3.0B to \$20.3B.

\subsection{Research Questions}

We address four questions:
\begin{enumerate}
\item \textbf{Single-Day Detection}: Can LLMs identify dealer hedging patterns when all temporal context is removed through obfuscation?
\item \textbf{Raw Chain Superiority}: Does strike-level data outperform parametric GEX summaries for structural detection?
\item \textbf{Regime Selectivity}: Can LLMs identify persistent 30-day regimes while rejecting transitional periods?
\item \textbf{Market Structure Evolution}: Did 0DTE proliferation create detectable structural change in dealer positioning regimes \citep{adams2025zerodte}?
\end{enumerate}

\subsection{Contributions}

This work makes four contributions: (i) \textbf{temporal obfuscation testing}, a generalizable methodology for distinguishing LLM structural reasoning from training-data memorization, validated through 2{,}221 evaluations across single-day and 30-day regime scales (1{,}412 real windows, 809 synthetic controls, 2020--2025); (ii) \textbf{raw-chain validation}, the first demonstration that LLMs detect dealer positioning more accurately from raw strike-level data than from pre-calculated GEX summaries (92.3\% vs.\ 61.5\%), establishing that scalar GEX is lossy compression of structural signal; (iii) \textbf{regime selectivity with negative controls}, with 69.1 percentage-point discrimination between 2024 persistent and 2020 fragmented markets and 0\% false positives on transitional and low-magnitude synthetic controls; and (iv) \textbf{detection-alpha orthogonality}, with detection stable at 68--74\% quarterly while economic profitability collapses (Sharpe 1.8 $\rightarrow$ 0.1), positioning the framework as a risk-management and surveillance tool rather than an alpha generator.

\subsection{Positioning}
\label{sec:introduction:positioning}

The contribution is primarily \textit{methodological}: temporal obfuscation testing, together with the WHO$\rightarrow$WHOM$\rightarrow$WHAT causal framework and the multi-scale validation protocol, is offered as a generalizable procedure for distinguishing LLM structural reasoning from training-data memorization. Options dealer gamma-exposure regime detection is the empirical demonstration domain because it combines mechanically grounded microstructure constraints, a large quantitative testbed (2{,}221 evaluations across six years), and a sharp pre- versus post-0DTE temporal contrast. The financial-market findings reported here---the 69.1 percentage-point 2024-vs-2020 detection gap, the 0\% false-positive rate on synthetic controls, and the gradual 2021--2024 regime evolution---are therefore presented as \textit{downstream evidence} that the methodology discriminates between persistent and fragmented market structures, not as novel claims about options market microstructure per se.

\subsection{Paper Organization}

Section~\ref{sec:related} reviews related work. Section~\ref{sec:methodology} presents the unified methodology covering obfuscation testing, causal framework, and regime detection criteria. Section~\ref{sec:single_day} reports single-day validation results including raw chain analysis. Section~\ref{sec:regime} presents multi-day regime detection and market structure evolution. Section~\ref{sec:discussion} discusses implications and limitations. Section~\ref{sec:conclusion} concludes.
